\section{Results III: Hardware and Platform Portability}

The first \YazSes{} evaluation was centered on one Linux laptop. The expanded campaign separates
two different questions that are often collapsed into the phrase cross-platform: whether the
same decoder produces similar text across operating-system and processor environments, and
whether the complete application path can be installed and exercised on those environments.

\subsection{Cross-platform decoding}

A common 60-utterance subset was decoded with CTranslate2 4.8.1 on Linux x86-64, Linux arm64,
macOS arm64, and Windows x86-64.

These are hosted CI runners, and operating system is confounded with processor: the four
runners reported an AMD EPYC 9V74, an Arm Neoverse-N2, a virtualized Apple M1, and an AMD EPYC
7763 respectively. The Linux x86-64 and Windows x86-64 columns therefore differ in CPU model as
well as in OS, so no column pair isolates an operating-system effect. What the experiment
measures is the joint environment a user actually gets.

% TODO(#508): generate from paper/results/platforms/*/wer.json.
\begin{table}[ht]
\centering
\caption{Cross-platform WER on the common 60-utterance subset.}
\label{tab:platform}
\begin{tabular}{lrrrrr}
\toprule
Checkpoint & Linux x86 & Linux arm64 & macOS arm64 & Windows x86 & Spread \\
\midrule
\texttt{tiny.en}  & 3.39 & 3.60 & 3.74 & 3.88 & 0.49 \\
\texttt{base.en}  & 3.32 & 3.32 & 3.25 & 3.39 & 0.14 \\
\texttt{small.en} & 2.05 & 2.05 & 2.05 & 2.05 & 0.00 \\
\bottomrule
\end{tabular}
\end{table}

The result is model-dependent rather than a single platform-wide effect. \texttt{small.en}
produces not only the same aggregate WER on all four runners but the same substitution,
deletion, insertion, and hit counts, whereas \texttt{tiny.en} shows a 0.49-point spread.
\texttt{base.en} is a reminder that an equal aggregate is not an equal output: the two Linux
runners both score 3.32\% from different error decompositions. The experiment therefore shows
that small aggregate WER differences need not reproduce identically across OS/ISA environments.
It does not establish that every checkpoint is meaningfully platform-sensitive, nor does it
identify which numerical kernel difference causes the observed spread.

\subsection{Installability is a different portability layer}

The companion dependency-resolution experiment simulates 92 target-platform/package-extra
combinations through the resolver, of which 84 resolve in the archived result. These are
dependency-resolution checks rather than 92 physical installations; failures expose packaging
constraints that are distinct from recognition accuracy.

These two experiments deliberately stop short of a stronger claim. A successful decoder run does
not exercise the entire interactive path: global activation, microphone permission, capture,
target focus, text injection, clipboard fallback, tray/daemon lifecycle, and accessibility APIs
remain separate OS-specific surfaces. End-to-end Windows and macOS validation is therefore
tracked as a separate experiment rather than inferred from Table~\ref{tab:platform}.